\section{Separated stages and episode bounds}\label{sec:stages}

Fix \(\beta=1/2\). Suppose \eqref{eq:5} fails along integers
\(N\to\infty\) for some fixed \(D_0\), and let \(Q_N\) be the
law of the environment conditioned on \(a_N<N^{-\beta}\) along this
sequence. Then
\[
 {\rm KL}(Q_N\mid P)\le D_0\log N.
 \tag{34}\label{eq:34}
\]
We will divide the interval of heights from zero to \(N\) into episodes.
During an episode, a retained law of a fixed number of particles is repeatedly extended:
a successful extension increases its separation scale, and a failed
extension decreases that scale. An episode ends when the scale drops below
a fixed floor. The construction below makes failures very unlikely under
the fresh environment law \(P\). We then use the displayed entropy bound
to control their expected number under \(Q_N\).

\subsection{Separated extensions}

\begin{proposition}[Separated stages]\label{stage-prop}
There are fixed constants \(f,g,\chi,K>0\) and an integer \(k\ge2\)
with
\[
 \frac1g<\frac{\beta}{64},\qquad
 \frac{(1+f/g)D_0}{K}<\frac{\beta}{64},
\]
having the following property. For every sufficiently large fixed \(b\),
put \(B=kb\) and choose a sufficiently large fixed floor
\(s_{\rm fl}\). At every level-count scale \(s\ge s_{\rm fl}\), put
\[
 s_-=se^{-fb},\qquad s_+=se^{gb},\qquad
 H_e=\lfloor se^{-b}\rfloor,\qquad H=\lfloor se^{\chi b}\rfloor.
\]
For every probability \(\pi\) on \(k\)-tuples supported on
\(\{{\rm sep}\ge r_s\}\), a fresh product environment satisfies
\[
\begin{array}{ll}
 \Gamma_h^\pi\{{\rm sep}(\text{endpoints})\ge r_{s_-}\}\ge e^{-B}
       &(1\le h\le H),\\
 \Gamma_H^\pi\{{\rm sep}(\text{endpoints})\ge r_{s_+}\}\ge e^{-B}
\end{array}
\tag{35}\label{eq:35}
\]
with probability at least \(1-e^{-2Kb}\). Each test uses the full
kernel from the same stage base; a passing test does not restrict the
paths used by later tests.
\end{proposition}

\begin{proof}
The main construction gives several chances to create separated mass
before any prescribed short interval of target heights. A failed chance
reveals only layers below the next available chance, so the estimate can
be repeated there. We first describe these chances and their probability
bounds, and then choose the constants to cover all target heights.
Throughout these estimates \(A\ge1\) and the integer \(k\ge2\)
are provisional fixed parameters. Write
\[
 \Lambda=\log16,\qquad j_*=\frac1{8A},\qquad
 m=\lfloor j_*b\rfloor.
\]
We take \(b\) large enough that \(m\ge1\). All uses of the rapid-loss
bound \eqref{eq:32} will have cap \(B_*=2kb\); after fixing \(b\),
increasing \(s_{\rm fl}\) will ensure its separation hypotheses.

\paragraph{A nested schedule before one target interval.}
Cover \([H_e,H]\) by integer cells \([v,v+w_m]\), where
\[
 w_0=16^m\left\lfloor\frac{H_e}{4\cdot16^m}\right\rfloor,
 \qquad w_i=16^{-i}w_0\quad(0\le i\le m),
\]
and \(v\) ranges through \(H_e,H_e+w_m,\ldots\) up to \(H\).
For the growth test at \(H\), add one cell with \(v=H\), using
\(H\) in place of \(H_e\) in the definition of \(w_0\).
After increasing \(s_{\rm fl}\), each \(w_0\) lies between one
eighth and one quarter of its corresponding horizon. A certificate for
the last cell may use paths above \(H\). These certificates will only
bound failure probabilities. The actual kernel tests in \eqref{eq:35}
alone determine the operational stopping layers.

Fix one cell. Its \(m\) opportunities have starting heights
\[
 v-2w_i,\qquad 0\le i<m.
\]
All these heights are nonnegative and increase with \(i\). When an
attempt begins at opportunity \(i\), use the normalized endpoint law of
the full original stream \(\Gamma_{v-2w_i}^{\pi}\). Uniform ellipticity makes this mass positive, and it uses only rows below \(v-2w_i\). In particular, after an
earlier failure we restart from this full stream, not from the mass
retained by the failed attempt.

The attempt first runs a seed kernel over \(w_i\) layers. Apply
\eqref{eq:30} with error \(e^{-6k}\), and put
\[
 \rho_i=r_{w_i/C_*}.
\]
Except with that error, it supplies mass at least \(g_0\) whose endpoint
separation is at least \(c_0\rho_i\), with \(c_0\le1\).
On failure we consume opportunity \(i\) and try \(i+1\), if one
remains. The failure is known at height \(v-w_i\), which is below
\(v-2w_{i+1}\).

On seed success, normalize its separated mass and protect it through the
whole target cell. Number the protection steps
\(l=0,\ldots,m-i-1\). Step zero starts at \(v-w_i\); a nonfinal
step \(l\) ends at \(v-2w_{i+l+1}\), and the final step ends at
\(v+w_m\). Every step has length at most
\(3w_i16^{-l}\). Thus, if failure is detected through a nonfinal
step \(l\), opportunity \(i+l+1\) begins exactly at its detection
height. The final step covers the entire cell and leaves no retry.


\paragraph{Retaining the separated seed.}
Choose an integer cutoff \(l_0\), sufficiently large as specified below.
Combine the first \(\min(l_0,m-i)\) protection steps into a single
kernel, ending at the same height as that last combined step. Its length
is at most \(3w_i=3C_*n(\rho_i)\).
Apply Lemma~\ref{kernel-initial-clipping} with
\[
 T=3C_*,\qquad \eta=c_0/4,\qquad p'=e^{-6kl_0}.
\]
Except with probability at most \(e^{-6kl_0}\), this retains mass
\(g_1>0\) with relative-motion budget \(c_0\rho_i/4\) throughout
the combined protection. On success normalize this mass. On failure
consume the \(\min(l_0,m-i)\) opportunities and restart at the next
opportunity if one remains. The test uses the kernel only through its
actual endpoint, including when fewer than \(l_0\) steps remain.

At every remaining step \(l\ge l_0\), allow relative-motion budget
\[
 d_l=(c_0/16)\rho_i2^{-l}.
\]
Let \(Z_l\) be minus the logarithm of the mass respecting this budget,
from the current normalized law. Continue the attempt if
\[
 \sum_{u=l_0}^{l}Z_u\le Ak(l+1),
\]
retaining and normalizing that budget-respecting mass. Otherwise declare
failure at the endpoint of this step. The early budget is
\(c_0\rho_i/4\), and all late budgets together are at most
\(c_0\rho_i/8\). Consequently all current laws have separation at
least \(c_0\rho_i/2\). Every restriction and normalization uses
only rows strictly below its endpoint.
Figure~\ref{fig:nested-opportunities} depicts the underlying height schedule.

\begin{figure}[tbp]
\centering
\begin{tikzpicture}[x=1cm,y=1cm,font=\small,>=Stealth]
  \fill[green!8] (10,-.4) rectangle (11.4,1.7);
  \node[anchor=south,text=green!40!black] at (10.7,1.7) {target cell};
  \draw[<->,green!45!black] (10,1.5) -- (11.4,1.5);
  \draw[->,black!55] (-.25,.9) -- (11.65,.9) node[right] {height};
  \foreach \x in {0,3,6,8.3,10,11.4}
    \draw[black!55] (\x,.8) -- (\x,1);
  \node[anchor=south] at (0,1.05) {$v-2w_i$};
  \node[anchor=south] at (3,1.05) {$v-w_i$};
  \node[anchor=south] at (6,1.05) {$v-2w_{i+1}$};
  \node[anchor=south] at (8.3,1.05) {$\cdots$};
  \node[anchor=north] at (10,.73) {$v$};
  \node[anchor=north] at (11.4,.73) {$v+w_m$};
  \draw[line width=3pt,orange!75!black] (0,0) -- (3,0);
  \node[anchor=south,text=orange!75!black] at (1.5,.1) {seed over $w_i$};
  \draw[line width=2pt,blue!65!black] (3,0) -- (11.4,0);
  \foreach \x in {3,6,8.3,9.25,11.4}
    \fill[blue!65!black] (\x,0) circle (2pt);
  \node[anchor=north,text=blue!65!black] at (4.5,-.08) {protection step $0$};
  \node[anchor=north,text=blue!65!black] at (8.8,-.08) {later protection steps};
  \draw[->,black!65] (6,-.5) -- (6,-1.05);
  \node[anchor=north,align=center,text width=11.8cm] at (5.7,-1.12)
    {Failure through step $l$: retry, if available, at\\
     $v-2w_{i+l+1}$, with index $i+l+1$.};
  \node[anchor=north,align=center,text width=11.8cm] at (5.7,-2.12)
    {Restart from the normalized \emph{full} stream
     $\Gamma_{v-2w_{i+l+1}}^{\pi}$, not from a failed retained substream.};
\end{tikzpicture}

\caption{The height schedule of one opportunity, not sample paths.
With $w_i=16^{-i}w_0$, seeding runs from $v-2w_i$ to $v-w_i$.
Nonfinal protection step $l$ ends at $v-2w_{i+l+1}$;
the final step ends at $v+w_m$ and covers the entire target cell.
A failed nonfinal protection consumes the traversed opportunities,
and any retry starts from the normalized full stream at its new base.
There is no retry after the final opportunity. The marked first
endpoint illustrates the geometric endpoint relation only. The initial
$\min(l_0,m-i)$ protection steps are tested together at their last
endpoint, so the diagram does not prescribe a retry at step zero.
Heights are not drawn to scale.}
\label{fig:nested-opportunities}
\end{figure}

\paragraph{Costs of the late protection steps.}
Set \(a'=2Ak\) and \(\lambda'=\lambda/4\). We will obtain,
conditionally at every active late step,
\[
\begin{aligned}
 E\big[e^{\lambda'(Z_l\wedge a'(l+1))}\mid\text{past}\big]
   &\le e^{C_{\rm mg}k},\\
 P(Z_l>a'(l+1)\mid\text{past})
   &\le e^{-(\lambda/2)a'(l+1)},
\end{aligned}
\tag{36}\label{eq:36}
\]
where \(C_{\rm mg}\) is independent of \(A\ge1\) and \(k\ge2\).
The cutoff \(l_0\), unlike \(C_{\rm mg}\), may depend on these
parameters, the seed constants, and fixed retry counts. The rapid-loss
estimate is uniform over normalized starting laws with the required
separation; it does not depend on the early retained mass \(g_1\).

Indeed, \eqref{eq:7} gives, for each fixed retry count \(L'\),
\[
 \frac{3w_i16^{-l}}{n(d_l/(2L'))}
 \le 3C_*\left(\frac{32L'}{c_0}\right)^2\left(\frac4{16}\right)^l
 =: C_{k,L'}e^{-\theta l},\qquad \theta=\log4.
\]
Choose a fixed \(c_x>0\) with
\((1+\lambda)c_x<\theta\) and \(c_x<1\). For
\(0\le x\le c_xl\), use \eqref{eq:32} with \(L'=1\).
After increasing \(l_0\), its second term is at most
\(e^{-\lambda x}\). Below the applicability threshold
\(k\log(1/\kappa)\), use the trivial bound one. Both ranges therefore
give
\[
 P(Z_l>x\mid\text{past})\le
 \min\{1,e^{C'k-\lambda x}\},\qquad 0\le x\le c_xl,
\]
with \(C'\) independent of \(A,k\).
For \(c_xl\le x\le a'(l+1)\), take a fixed integer
\(L'\) sufficiently large that
\(\theta L'>(1+2\lambda)a'\), and increase \(l_0\) again.
The same rapid-loss estimate then gives
\[
 P(Z_l>x\mid\text{past})
 \le 2^{L'}e^{CkL'-\lambda x}
       +(2CkC_{k,L'})^{L'}e^{x-\theta lL'}
 \le e^{-\lambda x/2}.
\]
Here \(l_0\) absorbs both the constants and the ellipticity threshold
for this fixed retry count. Also
\(a'(l+1)\le2Akm\le kb/4<B_*\).
Finally, integrating these tails with \(\lambda'=\lambda/4\) gives
\[
\begin{split}
 E[e^{\lambda'(Z_l\wedge a'(l+1))}\mid\text{past}]
 &\le 1+\lambda'\int_0^\infty e^{\lambda'x}
                  \min\{1,e^{C'k-\lambda x}\}\,dx
       +\lambda'\int_0^\infty e^{-\lambda x/4}\,dx\\
 &\le \frac43 e^{C'k/4}+1\le e^{C_{\rm mg}k}.
\end{split}
\]
This proves \eqref{eq:36} with the required uniformity.

\paragraph{Why all opportunities rarely fail.}
Now impose \(\lambda'A>C_{\rm mg}+8\). If the first cumulative
late failure occurs at \(l\), either \(Z_l>a'(l+1)\) or the
sum of the truncated costs exceeds \(Ak(l+1)\). Earlier passing
costs are below their own cutoffs, since they are nonnegative. Setting
inactive costs to zero, conditional iteration of \eqref{eq:36} bounds
the second event by
\[
 \exp\{C_{\rm mg}k(l-l_0+1)-\lambda'Ak(l+1)\}.
\]
Together with the cutoff tail, this is at most \(e^{-5k(l+1)}\).
A failed attempt consuming exactly \(n\) opportunities consequently
has probability at most \(e^{-2kn}\), measured at its start. The
possibilities are seed failure (\(n=1\)), early-protection failure
(\(n=\min(l_0,m-i)\)), and a first late failure (\(n=l+1\));
the stronger bounds just proved absorb any coincidence of these counts.

Each such failure is known before the next available opportunity, so
revealing the full stream up to the new start leaves its upper layers
fresh. Exhaustion partitions the \(m\) opportunities into positive
consumed blocks. For a specified composition of \(m\), conditional
iteration gives probability at most \(e^{-2km}\). There are
\(2^{m-1}\) compositions, so the probability that this cell has no
certificate is at most
\[
 e^{-(2k-\log2)m}.
\]

\paragraph{Choosing the constants.}
We can now choose constants in an order justified by the preceding
estimates. Fix \(C_{\rm mg}\), then take \(A\ge1\) with
\(\lambda'A>C_{\rm mg}+8\), and define
\[
 f=2+\Lambda j_*,\qquad \chi=g+1+\Lambda j_*.
\]
Choose \(g,K\) with
\(1/g<\beta/64\) and \((1+f/g)D_0/K<\beta/64\), and then choose
\(k\ge2\) so large that
\[
 \lambda k/4>2K+4,\qquad
 (2k-\log2)j_*>\chi+2+\Lambda j_*+2K+4.
 \tag{37}\label{eq:37}
\]
Next fix the seed constants and retry counts, choose \(l_0\) to
validate \eqref{eq:36}, and then obtain \(g_1\) from early clipping
with its specified error \(e^{-6kl_0}\). All these constants precede
\(b\). Thus even this very small early-protection error causes no
circularity: its possibly small retained mass \(g_1\) is fixed before
we enlarge \(b\). Finally enlarge \(s_{\rm fl}\), after \(b\), to
meet all scale and separation thresholds.

\paragraph{The full stream and the short heights.}
By \eqref{eq:31}, outside an event of probability at most
\(e^{-(2K+2)b}\) for sufficiently large \(b\),
\[
 \|\Gamma_h^\pi\|\ge q_\pi\ge e^{-B/4}\qquad\text{for every }h.
\]
We combine this exception with the cell failures by a union bound; none
of the cell-attempt estimates is conditioned on this full-stream event.

To cover heights up to \(H_e\), apply \eqref{eq:32} at horizon
\(H_e\), with \(d'=r_s/2\), \(x=B/2\), and a fixed retry count
\(L'>k/2+2K+4\). By \eqref{eq:7},
\[
 \frac{H_e}{n(d'/(2L'))}\le C(L')e^{-b}.
\]
Raising the rapid-loss bound to its indicated power gives an upper bound
of the form
\[
 C\bigl(e^{-\lambda kb/2}+e^{-(L'-k/2)b}\bigr)
 \le e^{-(2K+2)b}
\]
for large \(b\), by \eqref{eq:37}. Off this event, taking prefixes
gives mass at least \(e^{-B/2}\) at every height through \(H_e\)
with relative-motion budget \(r_s/2\). Initial separation is at least
\(r_s\), so these prefixes have separation at least
\(r_s/2\ge r_{s_-}\), again by \eqref{eq:7}.

\paragraph{Assembling the cell certificates.}
The total number of cells, including the separate growth cell, is at most
\[
 C\exp[(\chi+1+\Lambda j_*)b]+2
 \le\exp[(\chi+2+\Lambda j_*)b].
\]
On successful completion of the attempt from opportunity \(i\),
separation throughout that cell is at least \(c_0\rho_i/2\). The
scale inequality gives
\[
 n(c_0\rho_i/2)\ge\frac{c_0^2}{4C_*}w_i.
\]
For the traversal cells,
\[
 w_i\ge w_m\ge cse^{-(1+\Lambda j_*)b},
\]
which exceeds the required scale \(s_-=se^{-fb}\) after multiplication
by the preceding fixed constant. For the growth cell,
\[
 w_i\ge cse^{(\chi-\Lambda j_*)b},
\]
which similarly exceeds \(s_+=se^{gb}\). In each comparison the
remaining factor \(e^b\) absorbs the fixed constants for large \(b\).
Thus the certificate supplies the separation required in
\eqref{eq:35} at every height of its cell.

On the separately controlled event \(q_\pi\ge e^{-B/4}\), concatenate
the full stream to the attempt start with the restricted seed and
protection kernels. Their mass, and therefore their prefix mass at each
height of the cell, is at least
\[
 e^{-B/4}g_0g_1e^{-Akm}\ge e^{-B},
\]
because \(Akm\le B/8\) and \(g_0,g_1\) are fixed before \(b\).
The cell lies entirely in the protected part of the construction.

A union bound over cells, using \eqref{eq:37}, bounds the probability
of any missing cell certificate by \(e^{-(2K+2)b}\) for large \(b\).
Adding the full-stream and short-height exceptions proves
\eqref{eq:35}. At fixed \(b\), all the smallest scales tend uniformly
to infinity with \(s_{\rm fl}\). Enlarging that floor therefore
satisfies every scale and separation threshold used above.
\end{proof}

\subsection{Episode operation}

We now use the stage estimate to construct the actual episode boundaries.
Only the kernel tests in \eqref{eq:35} determine those boundaries; the
longer cell certificates in its proof do not enter the stopping rule.
Every retained measure below is a deterministic restriction of a kernel,
followed by normalization. No sample of a path is required.

\paragraph{Starting an episode.}
Start at height zero, and start a new episode after each reset below
\(N\). At an episode's initial height \(t<N\), take the global
normalized endpoint law
\[
 \mu_t=K_t(0,\cdot)/a_t
\]
in product \(k\) times. Inject this tuple law into layer \(t+1\)
with separation at least \(r_{s_{\rm fl}}\), using bounded elliptic
scripts. Concretely, choose \(k\) fixed positive integer offsets in
coordinate two, spaced by more than \(2r_{s_{\rm fl}}\). Prescribe
one particle's endpoint at a time, choosing the first offset whose
endpoint is at distance at least \(r_{s_{\rm fl}}\) in coordinate
two from every previously prescribed endpoint. Each earlier endpoint
excludes at most one of the \(k\) choices, so a choice always exists.
Move that particle laterally by its chosen offset and then directly up
one layer.

The scripts have a fixed length bound and hence cost at most a fixed
logarithmic mass \(B_0\). Normalize the retained tuple mass and set
\(s=s_{\rm fl}\). If the injection reaches \(N\), end the episode
there without a stage.

\paragraph{Testing one stage.}
At the current height \(u<N\), let \(\pi\) be the retained
normalized tuple law and let \(s\ge s_{\rm fl}\) be its scale.
Plan \(H=\lfloor se^{\chi b}\rfloor\) layers. Test the traversal
inequality in \eqref{eq:35} successively at offsets
\[
 h=1,\ldots,\min(H,N-u).
\]
Stop at the first failure. If offset \(H\) is reached, also test the
growth inequality there. All these tests use the same initial law
\(\pi\) and the full kernel from \(u\); passing an intermediate
test does not restrict the mass used for later tests.

\paragraph{Retaining mass and updating the scale.}
If the full planned stage passes, retain and normalize the growing
separated mass, and replace \(s\) by \(s_+=se^{gb}\).
If a test fails, replace \(s\) by \(s_-=se^{-fb}\) and retain mass
with the corresponding smaller separation. For a failed growth test,
the traversal test at the same height already supplies this mass. For
a first traversal failure at offset \(h\), extend the passing mass at
\(h-1\) by one direct upward step for every particle. At \(h=1\)
use the initial probability law as that passing mass. These upward
steps preserve separation, and in every failure case the retained mass
is at least
\[
 e^{-B}\kappa^k.
\]
Normalize it. If the updated scale is below \(s_{\rm fl}\), end the
episode and, if its endpoint is below \(N\), start the next episode
by the preceding injection rule. Otherwise continue with the next
stage at the new scale.

At height \(N\), end in all cases. A passing stage stopped at an
offset strictly less than \(H\) is called incomplete; its traversal
test supplies mass at least \(e^{-B}\). There is at most one such
stage. A failure encountered before this truncation still implies a
failure of the untruncated tests in \eqref{eq:35}, so it has the same
conditional probability bound under the fresh law \(P\).

Every test and retention uses only rows strictly below its actual
endpoint. Therefore all stage and episode boundaries are stopping times
for the layer filtration. In particular, the upper layers at every
subsequent stage start are fresh under \(P\).

Let \(J,F\) be the total numbers of stages and failed stages, and let
\(M\) be the number of episodes, with boundary heights
\(0=t_0<t_1<\cdots<t_M=N\). Put
\[
 h_i=t_{i+1}-t_i,\qquad
 J_i=\#\{\text{stages in episode }i\},\qquad
 c_i=\log(a_{t_i}/a_{t_{i+1}}),\qquad 0\le i<M.
\]

\begin{lemma}[Episode bounds]\label{stage-episodes}
For sufficiently large fixed \(b\), the episode construction satisfies
\(M\le J+1\) and
\[
 0\le c_i\le B_0/k+2bJ_i,\qquad
 \log(1+h_i)\le C(1+J_i).
 \tag{38}\label{eq:38}
\]
For large \(N\),
\[
 J\le\frac{\log N}{gb}+(1+f/g)F+1.
 \tag{39}\label{eq:39}
\]
Under \(Q_N\), fixed constants \(c,C>0\) and \(L_0<\infty\)
satisfy
\[
\begin{aligned}
 c\log N&\le E_{Q_N}M,& E_{Q_N}(M+J)&\le C\log N,\\
 E_{Q_N}\sum_{i<M}c_i\mathbf1_{\{J_i\le L_0\}}
   &\ge c\log N.&&
\end{aligned}
\tag{40}\label{eq:40}
\]
With the padding \(t_i=N\) for \(i\ge M\), each \(t_i\) is a
bounded layer stopping time, and \(\{i<M\}\) is known at \(t_i\).
\end{lemma}

\begin{proof}
\smallskip\noindent\textbf{Mass and height within an episode.}
The unrestricted quenched \(k\)-survival mass through \(t_i\) is
\(a_{t_i}^k\). Starting from its normalized endpoint law, our retained
continuation costs at most
\(B_0+(kb+k\log(1/\kappa))J_i\) in logarithmic mass. Its
concatenation is contained in the original survival kernel through
\(t_{i+1}\), so
\[
 a_{t_{i+1}}^k\ge a_{t_i}^k
 \exp\{-B_0-(kb+k\log(1/\kappa))J_i\}.
\]
This proves the upper bound for \(c_i\) in \eqref{eq:38} once
\(b\ge\log(1/\kappa)\); its nonnegativity follows from the
monotonicity of \(a_t\). Each stage length is at most
\(se^{\chi b}\), and the scale increases from the fixed floor by
at most a factor \(e^{gb}\) per stage. Summing these geometric
bounds, together with the one-layer injection, proves the height bound
in \eqref{eq:38}. Every episode except possibly the last contains a
failed stage, so \(M\le J+1\). The last episode may consist only of
its injection.

\smallskip\noindent\textbf{Counting stages by their scale changes.}
For this accounting only, update the scale after the last completed or
failed stage as if the operation continued, and clamp it upward to the
floor after a reset. Every completed success consumes its full planned
height \(H\ge s_+\), so its updated scale is at most
\(\max(s_{\rm fl},N)\). Failures decrease the scale, and the
clamping returns it only to \(s_{\rm fl}\). If \(S\) is the
number of completed successes, their increases of \(gb\), minus the
losses of at most \(fb\) per failure, therefore satisfy
\[
 gbS-fbF\le\log(N/s_{\rm fl})\le\log N
\]
for large \(N\). There is at most one passing incomplete stage, so
\(J\le S+F+1\). This proves \eqref{eq:39}.

\smallskip\noindent\textbf{Transferring the failure bound to \(Q_N\).}
Under \(P\), conditional on the past at any actual stage start, its
failure probability is at most \(e^{-2Kb}\). If \(I\) is its
failure indicator, then
\[
\begin{split}
 E_P[e^{KbI-e^{-Kb}}\mid\text{stage past}]
 &\le e^{-e^{-Kb}}
       \bigl(1+e^{-2Kb}(e^{Kb}-1)\bigr)\le1.
\end{split}
\]
There are at most \(N\) actual stages, since every stage advances at
least one layer. Pad the list to \(N\) slots, with active indicator
\(A_r\), and put \(I_r=0\) in inactive slots. Before each slot its
activity is known, and
\[
 E_P\big[e^{KbI_r-e^{-Kb}A_r}\mid\text{past}\big]\le1.
\]
The inactive factor is one. Conditional iteration, using
\(\sum_r I_r=F\) and \(\sum_r A_r=J\), gives
\[
 E_P e^{KbF-e^{-Kb}J}\le1.
\]
Applying the entropy inequality with \eqref{eq:34} yields
\[
 KbE_{Q_N}F\le D_0\log N+e^{-Kb}E_{Q_N}J.
\]
Only the exponential estimate was taken under the fresh law; this last
inequality requires no independence under \(Q_N\).

Write \(c_f=1+f/g\). Combining the last inequality with
\eqref{eq:39} gives
\[
 E_{Q_N}J\le
 \frac{\log N}{b}\left(\frac1g+\frac{c_fD_0}{K}\right)
 +\frac{c_fe^{-Kb}}{Kb}E_{Q_N}J+1.
\]
The coefficient in parentheses is less than \(\beta/32\). After
absorbing the exponentially small term, for large fixed \(b\) and
then large \(N\),
\[
 2bE_{Q_N}J\le\frac\beta4\log N.
\]

\smallskip\noindent\textbf{Positive loss in episodes with bounded stage count.}
Under \(Q_N\),
\(\sum_{i<M}c_i=-\log a_N>\beta\log N\). Summing
\eqref{eq:38} and using the last bound shows
\[
 \frac{B_0}{k}E_{Q_N}M
 \ge\beta\log N-2bE_{Q_N}J
 \ge\frac{3\beta}{4}\log N.
\]
This proves the lower bound on \(E_{Q_N}M\); its upper bound, together
with that on \(E_{Q_N}J\), follows from \(M\le J+1\).
For any \(L_0>0\), the loss in the remaining episodes satisfies
\[
 \sum_{i<M}c_i\mathbf1_{\{J_i>L_0\}}
 \le\left(\frac{B_0}{kL_0}+2b\right)J.
\]
Although the injection cost \(B_0\) may be large, it is fixed. Choose
\(L_0\) after it so that the first term's expected contribution is
at most \((\beta/4)\log N\). Episodes with more than \(L_0\)
stages then account for at most \((\beta/2)\log N\) of expected
loss. Subtracting this from the total proves the last assertion of
\eqref{eq:40}.

Finally, the operational tests use only rows below their tested layer,
so all episode boundaries are bounded layer stopping times. With the
stated padding, \(\{i<M\}=\{t_i<N\}\) is known at \(t_i\).
\end{proof}
